\hypertarget{appendix-0659864c0aa3}{}\label{referencestart:appendix-0659864c0aa3}
\section{Vectors coordinates and affine points}\label{appendix-0659864c0aa3-vectors-coordinates-and-affine-points}

\hypertarget{appendix-0659864c0aa3-vectors-coordinates-and-affine-points}{}

\textbf{Curriculum depth:} Introduced · \textbf{First seen:} Phase6 · \textbf{Latest development:} Phase6

\begin{allmightroledefinition}

\subsection{Definition and mechanism}\label{appendix-0659864c0aa3-definition-and-mechanism}

\hypertarget{appendix-0659864c0aa3-definition-and-mechanism}{}

A free vector describes movement independently of its starting point. A coordinate tuple describes that vector relative to an ordered basis. An affine point is a location without an intrinsic preferred origin. Point differences produce vectors; point translations produce points.

\end{allmightroledefinition}

For the workcell panel, \(B-A=(4,0,0)\) is a displacement in metres. Changing the origin changes point coordinates but cancels in this difference.

\subsection{Prerequisites}\label{appendix-0659864c0aa3-prerequisites}

\hypertarget{appendix-0659864c0aa3-prerequisites}{}

\begin{itemize}
\tightlist
\item
  \hyperref[appendix-52051ad550fe]{Scalars fields and real-number models} --- Field operations define which scalar multipliers the vector laws may use.
\end{itemize}

\subsection{Used by}\label{appendix-0659864c0aa3-used-by}

\hypertarget{appendix-0659864c0aa3-used-by}{}

\begin{itemize}
\tightlist
\item
  \hyperref[appendix-173b52c063ae]{Vector spaces subspaces and span} --- Builds on this mechanism in the connected geometry story.
\item
  \href{../04-Light-Transport-and-Image-Synthesis/main.pdf\#nameddest=appendix-02b9d9a81dca}{Barycentric coordinates and interpolation (Light Transport and Image Synthesis, Reference notes: Barycentric coordinates and interpolation)} --- Builds on this mechanism in the connected geometry story.
\item
  \href{../04-Light-Transport-and-Image-Synthesis/main.pdf\#nameddest=appendix-cb3b26862282}{Homogeneous coordinates and perspective division (Light Transport and Image Synthesis, Reference notes: Homogeneous coordinates and perspective division)} --- Builds on this mechanism in the connected geometry story.
\item
  \hyperref[appendix-37a39984ba66]{Vector value types and component access} --- Builds on this mechanism in the connected geometry story.
\item
  \hyperref[appendix-380a220e2529]{Coordinate-space types and validated invariants} --- Builds on this mechanism in the connected geometry story.
\end{itemize}

\subsection{Recall and next depth}\label{appendix-0659864c0aa3-recall-and-next-depth}

\hypertarget{appendix-0659864c0aa3-recall-and-next-depth}{}

Show why point plus point needs an affine-weight policy, while point plus displacement does not. Full production proofs, target-specific integration and learner mastery require further evidence.

\subsection{Source and navigation}\label{appendix-0659864c0aa3-source-and-navigation}

\hypertarget{appendix-0659864c0aa3-source-and-navigation}{}

\href{../02-Mathematics-for-Visual-Computing/main.pdf\#nameddest=01-foundations-i-ch017-an-arrow-is-not-the-address-written-beside-it}{Teaching narrative (Mathematics for Visual Computing, Ch. 1)} · \hyperref[appendix-fd090d982a15-phase-6---parts-501-600]{Coverage and source gap} · \hyperref[appendix-a103e7846d21-phase-6---source-analysis-and-verification]{Evidence and limits} · \hyperref[appendix-4ad47bdd269a]{Dependency Map}
